\documentclass[11pt]{beamer}

\usetheme{Madrid}
\usecolortheme{default}
\usepackage{graphicx}
\usepackage{hyperref}
\usepackage{tikz}
\usepackage{qrcode}

\setbeamertemplate{footline}[frame number]

\title{Formalator}
\subtitle{Bidirectional Text Translator}
\author{Amir Gafarov}
\institute{
    a.gafarov@innopolis.university \\
    CSE-03
}
\date{\today}

\begin{document}

% ===== TITLE SLIDE =====
\begin{frame}
    \titlepage
\end{frame}

% ===== CONTEXT =====
\begin{frame}{Context}

    \vspace{0.3cm}
    \textbf{End-User}
    \begin{itemize}
        \item Anyone who needs to adapt their text tone — from students writing professional emails to professionals who want a more casual, approachable tone.
    \end{itemize}

    \vspace{0.3cm}
    \textbf{Problem}
    \begin{itemize}
        \item Many people struggle to adjust their writing style for different contexts (formal vs.\ informal). Getting it right takes time and effort.
    \end{itemize}

    \vspace{0.3cm}
    \textbf{One-Sentence Pitch}
    \begin{itemize}
        \item \textit{Formalator} is a web app and Telegram bot that instantly rewrites your text in a formal or informal style using LLM — two interfaces, one engine, both directions.
    \end{itemize}

\end{frame}

% ===== IMPLEMENTATION =====
\begin{frame}{Implementation}

    \textbf{Architecture}
    \begin{itemize}
        \item \textbf{Backend} — FastAPI (Python) with LLM client
        \item \textbf{Frontend} — React + TypeScript + Vite + Bootstrap 5
        \item \textbf{Telegram Bot} — aiogram 3.x (Python)
        \item \textbf{LLM} — OpenAI-compatible API
        \item \textbf{Deployment} — Docker Compose (3 services + nginx)
    \end{itemize}

    \vspace{0.3cm}
    \textbf{Version 1} — Formal$\,\rightarrow\,$Informal only, Telegram bot, CLI testing.

    \vspace{0.3cm}
    \textbf{Version 2} — Bidirectional translation, React web UI, \texttt{/mode} bot command, unified \texttt{direction} parameter in API.

\end{frame}

% ===== IMPLEMENTATION — FEEDBACK =====
\begin{frame}{Implementation — Feedback Addressed}

    \textbf{TA Feedback Points Resolved}
    \begin{itemize}
        \item \textbf{Docker build consistency} — Fixed \texttt{npm ci} crash on VM by switching to \texttt{npm install}; frontend now uses the same package manager (\texttt{pnpm}) locally and in CI.
        \item \textbf{Multi-interface parity} — Both web UI and Telegram bot support the same bidirectional translation feature.
        \item \textbf{System prompt quality} — Separate prompts for each direction with clear style rules (contractions, vocabulary, tone).
        \item \textbf{Deployment reliability} — Docker Compose with proper \texttt{.dockerignore}, build context separation, and \texttt{--no-cache} rebuild on VM.
    \end{itemize}

\end{frame}

% ===== DEMO =====
\begin{frame}{Demo}

    \begin{center}
        \vspace{1cm}
        {\Huge 🎬}

        \vspace{0.5cm}
        \textbf{\Large Video Demonstration — Version 2}

        \vspace{0.5cm}
        \footnotesize Pre-recorded demo with voice-over (≤2 min)

        \vspace{0.3cm}
        \textit{Shows:}
        \begin{itemize}
            \item Web UI: bidirectional translation, direction toggle, examples
            \item Telegram Bot: \texttt{/mode} switch, live translation
        \end{itemize}

        \vspace{0.5cm}
        {\small \textit{(Insert video here)}}
    \end{center}

\end{frame}

% ===== LINKS =====
\begin{frame}{Links}

    \vspace{0.2cm}
    \begin{columns}[T]
        \begin{column}{0.48\textwidth}
            \textbf{GitHub Repository}
            \vspace{0.2cm}
            \begin{center}
                \scriptsize
                \url{https://github.com/Omarichev/Formalator}
                \vspace{0.2cm}

                \qrcode[height=2.5cm]{https://github.com/Omarichev/Formalator}
            \end{center}
        \end{column}

        \begin{column}{0.48\textwidth}
            \textbf{Deployed Product}
            \vspace{0.2cm}
            \begin{center}
                \scriptsize
                \url{http://<your-vm-ip>:80}
                \vspace{0.2cm}

                \qrcode[height=2.5cm]{http://<your-vm-ip>:80}
            \end{center}
        \end{column}
    \end{columns}

\end{frame}

\end{document}
